% Options for packages loaded elsewhere
\PassOptionsToPackage{unicode}{hyperref}
\PassOptionsToPackage{hyphens}{url}
\documentclass[
]{article}
\usepackage{xcolor}
\usepackage{amsmath,amssymb}
\setcounter{secnumdepth}{5}
\usepackage{iftex}
\ifPDFTeX
  \usepackage[T1]{fontenc}
  \usepackage[utf8]{inputenc}
  \usepackage{textcomp} % provide euro and other symbols
\else % if luatex or xetex
  \usepackage{unicode-math} % this also loads fontspec
  \defaultfontfeatures{Scale=MatchLowercase}
  \defaultfontfeatures[\rmfamily]{Ligatures=TeX,Scale=1}
\fi
% \usepackage{lmodern}
\ifPDFTeX\else
  % xetex/luatex font selection
\fi
% Use upquote if available, for straight quotes in verbatim environments
\IfFileExists{upquote.sty}{\usepackage{upquote}}{}
\IfFileExists{microtype.sty}{% use microtype if available
  \usepackage[]{microtype}
  \UseMicrotypeSet[protrusion]{basicmath} % disable protrusion for tt fonts
}{}
\makeatletter
\@ifundefined{KOMAClassName}{% if non-KOMA class
  \IfFileExists{parskip.sty}{%
    \usepackage{parskip}
  }{% else
    \setlength{\parindent}{0pt}
    \setlength{\parskip}{6pt plus 2pt minus 1pt}}
}{% if KOMA class
  \KOMAoptions{parskip=half}}
\makeatother
\usepackage{longtable,booktabs,array}
\newcounter{none} % for unnumbered tables
\usepackage{calc} % for calculating minipage widths
% Correct order of tables after \paragraph or \subparagraph
\usepackage{etoolbox}
\makeatletter
\patchcmd\longtable{\par}{\if@noskipsec\mbox{}\fi\par}{}{}
\makeatother
% Allow footnotes in longtable head/foot
\IfFileExists{footnotehyper.sty}{\usepackage{footnotehyper}}{\usepackage{footnote}}
\makesavenoteenv{longtable}
\setlength{\emergencystretch}{3em} % prevent overfull lines
\providecommand{\tightlist}{%
  \setlength{\itemsep}{0pt}\setlength{\parskip}{0pt}}
\usepackage[]{natbib}
\bibliographystyle{plainnat}
\usepackage{bookmark}
\IfFileExists{xurl.sty}{\usepackage{xurl}}{} % add URL line breaks if available
\urlstyle{same}
\hypersetup{
  hidelinks,
  pdfcreator={LaTeX via pandoc}}

\author{}
\date{}

\title{Cognitive Integrity Framework: A Qualitative Review\\\normalsize Part 3 of 3: Synthesizing Theory and Computation for Practitioners}
\author{Daniel Ari Friedman\\\footnotesize{Active Inference Institute}\\\footnotesize{\texttt{daniel@activeinference.institute}}\\\footnotesize{\href{https://orcid.org/0000-0001-6232-9096}{ORCID: 0000-0001-6232-9096}} \\ \footnotesize{2026-02-15}}

% Core mathematical packages
\usepackage{amsmath,amssymb,amsthm}
\usepackage{mathtools}

% Algorithm formatting
\usepackage{algorithm}
\usepackage{algpseudocode}

% Tables
\usepackage{booktabs}
\usepackage{multirow}
\usepackage{array}
\usepackage{longtable}

% Graphics
\usepackage{graphicx}
\usepackage{tikz}
\usepackage{pgfplots}
\usetikzlibrary{shapes,arrows,positioning,calc,fit,backgrounds}
\pgfplotsset{compat=1.18}

% Cross-referencing with smart naming
\usepackage{hyperref}
\usepackage[capitalise,noabbrev,nameinlink]{cleveref}

% Configure cleveref for custom environments
\crefname{definition}{Definition}{Definitions}
\crefname{theorem}{Theorem}{Theorems}
\crefname{lemma}{Lemma}{Lemmas}
\crefname{property}{Property}{Properties}
\crefname{corollary}{Corollary}{Corollaries}
\crefname{algorithm}{Algorithm}{Algorithms}
\crefname{equation}{Equation}{Equations}
\crefname{table}{Table}{Tables}
\crefname{figure}{Figure}{Figures}

% Theorem environments with consistent numbering
\newtheorem{definition}{Definition}[section]
\newtheorem{theorem}{Theorem}[section]
\newtheorem{lemma}[theorem]{Lemma}
\newtheorem{corollary}[theorem]{Corollary}
\newtheorem{property}{Property}[section]
\newtheorem{axiom}{Axiom}[section]

% Math operators
\DeclareMathOperator*{\argmax}{arg\,max}
\DeclareMathOperator*{\argmin}{arg\,min}
\DeclareMathOperator{\sign}{sign}
\DeclareMathOperator{\KL}{KL}
\DeclareMathOperator{\Tr}{Tr}

% Custom commands for notation consistency
\newcommand{\calA}{\mathcal{A}}
\newcommand{\calB}{\mathcal{B}}
\newcommand{\calC}{\mathcal{C}}
\newcommand{\calD}{\mathcal{D}}
\newcommand{\calF}{\mathcal{F}}
\newcommand{\calG}{\mathcal{G}}
\newcommand{\calH}{\mathcal{H}}
\newcommand{\calI}{\mathcal{I}}
\newcommand{\calM}{\mathcal{M}}
\newcommand{\calO}{\mathcal{O}}
\newcommand{\calP}{\mathcal{P}}
\newcommand{\calS}{\mathcal{S}}
\newcommand{\calT}{\mathcal{T}}
\newcommand{\calW}{\mathcal{W}}
% Note: \Phi is a standard LaTeX command, do not redefine

% Trust notation
\newcommand{\trust}[2]{\mathcal{T}_{#1 \to #2}}
\newcommand{\trustt}[3]{\mathcal{T}_{#1 \to #2}^{#3}}

% Belief notation
\newcommand{\belief}[2]{\mathcal{B}_{#1}(#2)}
\newcommand{\belieft}[3]{\mathcal{B}_{#1}^{#2}(#3)}

% Cognitive state
\newcommand{\cogstate}[1]{\sigma_{#1}}

% Attack notation
\newcommand{\attack}[1]{\mathcal{A}_{#1}}
\newcommand{\adversary}[1]{\Omega_{#1}}

% Defense notation
\newcommand{\firewall}{\mathcal{F}}
\newcommand{\sandbox}{\mathcal{S}_{box}}

% Probability and expectation
\newcommand{\E}{\mathbb{E}}
\newcommand{\Prob}{\mathbb{P}}
\newcommand{\indicator}{\mathbb{1}}

% QED symbol
\renewcommand{\qedsymbol}{$\blacksquare$}

% Page Layout (Slightly smaller margins)
\usepackage[margin=1in]{geometry}

% Hyperlink Styling
\hypersetup{
    colorlinks=true,
    linkcolor=red,
    filecolor=red,
    urlcolor=red,
    citecolor=red
}

\begin{document}

\maketitle
\thispagestyle{empty}


\vspace*{2cm}

\begin{center}
\begin{minipage}{0.7\textwidth}
\centering
\Large\itshape
``In theory, there is no difference between\\[0.3em]
theory and practice. In practice, there is.''
\vspace{1em}

\normalsize\upshape
--- Yogi Berra (attributed)
\end{minipage}
\end{center}

\vspace{2cm}

\section{Abstract}\label{abstract}

Multiagent AI systems---autonomous coding assistants, research
pipelines, financial decision engines---have moved from prototype to
production in under two years. With them comes a new class of security
concern: attacks that target not data or infrastructure but the
\emph{reasoning processes} of AI agents. Prompt injections that
propagate through delegation chains, trust relationships that launder
adversarial influence, and coordination mechanisms vulnerable to
strategic manipulation all represent cognitive attack surfaces absent
from traditional security models.

The \textbf{Cognitive Integrity Framework (CIF)}, presented across two
companion papers, offers the first comprehensive formal and
computational treatment of this problem. Part 1 establishes mathematical
foundations: a trust calculus with provably bounded delegation, defense
composition algebras with multiplicative detection guarantees, and
information-theoretic limits on attack stealth. Part 2 provides
computational validation: eight implemented defense modules (1,594
passing tests), a 950-attack corpus spanning four threat categories, and
parametric architecture-aware simulation across four production
multiagent topologies.

This paper is a qualitative review and practitioner's guide to the CIF
series. We synthesize the key insights from both papers into accessible
language, contextualize the formal results within the current deployment
landscape, assess what the research has established and where gaps
remain, and distill practical recommendations for teams building and
operating multiagent AI systems. No formal prerequisites are assumed;
readers seeking mathematical detail are referred to Parts 1 and 2.

\subsection{Paper Series}\label{paper-series}

\textbf{DOI}: 10.5281/zenodo.18364130

This is Part 3 of the \emph{Cognitive Security for Multiagent Operators}
series:

\begin{itemize}
\tightlist
\item
  \textbf{Part 1} (DOI: 10.5281/zenodo.18364119): Formal foundations and
  theoretical analysis
\item
  \textbf{Part 2} (DOI: 10.5281/zenodo.18364128): Computational
  validation and implementation
\item
  \textbf{Part 3} (this paper): Qualitative review and practitioner's
  synthesis
\end{itemize}

\newpage

\section{Why Cognitive Security Matters Now}\label{sec:introduction}

\subsection{The Operational Reality}\label{the-operational-reality}

Something fundamental changed in how AI systems work, and the security
community is catching up.

In 2023, AI security often meant preventing chatbots from saying things
they shouldn't. The attack surface was a text box; the defense was a
filter.

By 2026, we are securing \textbf{multiagent operators}---networks of
specialized AI agents that delegate to each other, form beliefs about
each other's outputs, build trust relationships over time, and take
actions with real-world consequences. These systems write code, manage
infrastructure, and move money.

The shift is from ``content safety'' to ``cognitive integrity.'' The
risk isn't just that an agent says something wrong, but that it
\emph{believes} something wrong---and acts on it.

\subsection{The Good News: It's
Solvable}\label{the-good-news-its-solvable}

This is not a theoretical warning about future doom. It is an
engineering problem with established solutions.

The Cognitive Integrity Framework (CIF) was developed to secure these
systems, and the first two papers in this series demonstrated its
efficacy.

\begin{itemize}
\tightlist
\item
  \textbf{Part 1: Formal Foundations} proved that trust can be
  mathematically bounded. We defined the ``Trust Calculus'' which
  guarantees that no matter how clever an adversary is, they cannot
  amplify their influence through delegation chains.
\item
  \textbf{Part 2: Computational Validation} implemented this theory in
  Python and tested it against a corpus of 950 attacks across four
  production architectures.
\end{itemize}

The result was \textbf{1,594 passing tests and a 97\% structural
constraint satisfaction rate} against direct injection attacks in fully
defended configurations.

\subsection{The Purpose of This Guide}\label{the-purpose-of-this-guide}

We wrote Part 1 for the theorists and Part 2 for the experimentalists.
We wrote this paper---Part 3---to translate those findings into
practice.

Our goal is to describe how the defenses validated in the previous
papers can be architected in production systems. We focus on the
practical application of the formal proofs:

\begin{itemize}
\tightlist
\item
  How the \textbf{Trust Decay} factor (\(\delta\)) functions in
  different topologies.
\item
  How \textbf{Behavioral Tripwires} served as effective detection
  mechanisms for hallucination.
\item
  How the \textbf{Cognitive Firewall} filtered inputs before they became
  beliefs.
\end{itemize}

\subsection{How to Use This Resource}\label{how-to-use-this-resource}

\begin{itemize}
\tightlist
\item
  \textbf{Section 2} summarizes the theoretical concepts from Part 1,
  providing the necessary vocabulary.
\item
  \textbf{Section 3} reviews the empirical evidence from Part 2,
  detailing which architectures performed best against specific threats.
\item
  \textbf{Section 4} analyzes the attack scenarios used in our testing
  corpus.
\item
  \textbf{Section 5} presents the specific configuration profiles that
  yielded the highest security margins in simulation.
\item
  \textbf{Sections 6-7} discuss the limitations discovered during
  testing and the open problems that remain.
\end{itemize}

This paper serves as a report on the current state of cognitive security
engineering, grounded in the data and definitions of the CIF series.

\newpage

\section{The Formal Foundation: Concepts from Paper
1}\label{sec:theory-review}

Part 3 builds directly on the formal framework established in Part 1.
For clarity, we summarize the core definitions and theorems here,
utilizing the notation defined in the formal manuscript.

\subsection{\texorpdfstring{1. The Adversary Hierarchy
(\(\Omega\))}{1. The Adversary Hierarchy (\textbackslash Omega)}}\label{the-adversary-hierarchy-omega}

Paper 1 formalized the ``Scope of Threat'' through a hierarchical
taxonomy. This hierarchy allows precise definition of defensive scope.

\begin{itemize}
\tightlist
\item
  \textbf{\(\Omega_1\) (External)}: The adversary controls inputs (e.g.,
  prompt injection). The agent's internal state is intact.
\item
  \textbf{\(\Omega_2\) (Peripheral)}: The adversary controls tools or
  RAG data (e.g., poisoned retrieval). The agent's perception is
  compromised.
\item
  \textbf{\(\Omega_3\) (Agent)}: The adversary controls the agent's
  weights or context (e.g., identity implementation). The agent itself
  is untrusted.
\item
  \textbf{\(\Omega_4\) (Coordination)}: The adversary controls a subset
  of the swarm (e.g., Sybil agents). The consensus mechanism is under
  attack.
\item
  \textbf{\(\Omega_5\) (Systemic)}: The adversary controls the
  orchestrator or infrastructure. The system's rules are compromised.
\end{itemize}

The simulations in Part 2 demonstrated that defense difficulty scales
non-linearly with this hierarchy. While \(\Omega_1\) attacks were
consistently blocked by surface-level filters (96\%+), \(\Omega_4\)
attacks required coordination-level protocols to detect (74\%).

\subsection{\texorpdfstring{2. The Trust Calculus
(\(T\))}{2. The Trust Calculus (T)}}\label{the-trust-calculus-t}

A central contribution of Paper 1 is the \textbf{Trust Calculus}, a
formal system for reasoning about belief reliability. It defines Trust
(\(T\)) not as a binary permission but as a continuous property of a
belief \(b\), denoted as \(T(b) \in [0, 1]\).

The formal definition of Trust Update (Theorem 3.1 in Paper 1)
establishes that trust must decay across delegation chains:

\begin{quote}
\textbf{Theorem 3.1 (Trust Preservation)}: \emph{For any delegation
chain \(C = \{a_1 \to a_2 \to \dots \to a_n\}\), the trust in the final
output cannot exceed the trust of the weakest link, degraded by the
distance from the source.}
\[ T(result) \le \min_{i} T(a_i) \cdot \delta^{|C|} \] \emph{where
\(\delta\) is the decay factor (\(0 < \delta < 1\)).}
\end{quote}

This theorem provides the mathematical basis for the ``Trust Decay''
mechanism evaluated in Part 2. It ensures that uncertainty is preserved
and amplified effectively as information travels through the network.

\subsection{\texorpdfstring{3. The Cognitive Firewall
(\(\Phi\))}{3. The Cognitive Firewall (\textbackslash Phi)}}\label{the-cognitive-firewall-phi}

The \textbf{Cognitive Firewall} is defined in Paper 1 as a function
\(\Phi\) that maps inputs to decisions based on three verification
layers:

\begin{enumerate}
\def\labelenumi{\arabic{enumi}.}
\tightlist
\item
  \textbf{Syntactic Verification (\(V_{syn}\))}: Checks for structural
  anomalies.
\item
  \textbf{Semantic Verification (\(V_{sem}\))}: Checks for meaning-level
  violations.
\item
  \textbf{Pragmatic Verification (\(V_{prag}\))}: Checks for contextual
  anomalies.
\end{enumerate}

In the Part 2 experiments, this modular structure was shown to be the
primary defense against \(\Omega_1\) (External) attacks.

\subsection{4. The Stealth-Impact
Tradeoff}\label{the-stealth-impact-tradeoff}

Paper 1 provides a theoretical bound on attack performance, formalized
as the Stealth-Impact Tradeoff.

\begin{quote}
\textbf{Stealth-Impact Tradeoff}: \emph{For a given defense sensitivity
\(\epsilon\), the probability of detection \(P(d)\) approaches 1 as the
divergence of the attack behavior from the baseline increases.}
\end{quote}

This formalism suggests that catastrophic attacks are inherently easier
to detect than subtle attacks. Part 2's data consistently validated
this: High-impact attacks were detected 98\% of the time, while
low-impact attacks were detected only 74\% of the time.

\subsection{5. Defense Composition}\label{defense-composition}

Finally, Paper 1 defines the \textbf{Composition Algebra}, determining
how output probabilities of distinct modules interact. The key result is
that orthogonal defenses compose multiplicatively.

This ``Swiss Cheese Model'' was empirically validated in Part 2, where
the full stack (97\% structural constraint satisfaction) significantly
outperformed the sum of its parts.

\newpage

\section{The Evidence: What We Proved in Part
2}\label{sec:paper2-review}

Part 1 was the theory. Part 2 was the proof. We ran \textbf{1,594 tests}
against \textbf{950 attack samples} across \textbf{four production
architectures}.

Here is what the data says.

\subsection{The Experimental Setup}\label{the-experimental-setup}

We tested four common multi-agent architectures found in production:

\begin{itemize}
\tightlist
\item
  \textbf{Claude Code} (Hierarchical)
\item
  \textbf{AutoGPT} (Autonomous Loop)
\item
  \textbf{CrewAI} (Role-Based Team)
\item
  \textbf{LangGraph} (State Machine)
\end{itemize}

The test corpus included direct prompt injection, poisoned RAG contexts,
deep trust exploitation, and multi-turn social engineering.

\subsection{Finding 1: Defense Layering vs.~Individual
Efficacy}\label{finding-1-defense-layering-vs.-individual-efficacy}

\textbf{The Data}: Individual defenses (like just a firewall) stopped
\textasciitilde60-70\% of attacks. The full CIF stack stopped
\textbf{97\%}. \textbf{The Implication}: The defenses demonstrated
orthogonal coverage. The firewall blocked inputs that the sandbox would
have missed, and the sandbox identified anomalies that the trust
calculus would have permitted. The data suggests that removing any
single layer creates a statistically significant vulnerability gap.

\subsection{Finding 2: State Machine
Determinism}\label{finding-2-state-machine-determinism}

\textbf{The Data}: \textbf{LangGraph} architectures achieved the highest
detection rates (98\%). \textbf{The Mechanism}: In a state machine,
valid transitions are explicitly defined. Attempts to move from an
\texttt{Analyze} state to an \texttt{Execute} state without passing
through \texttt{Verify} were caught immediately. \textbf{The
Implication}: Explicit state definitions provide a ``security for free''
effect, where the architecture itself enforces behavioral invariants
that are difficult to bypass.

\subsection{Finding 3: Hierarchical
Vulnerability}\label{finding-3-hierarchical-vulnerability}

\textbf{The Data}: \textbf{Claude Code} and \textbf{AutoGPT} styles were
strong against external attacks but vulnerable to \textbf{Systemic
(\(\Omega_5\))} compromise. \textbf{The Mechanism}: When the ``Boss''
agent was tricked, it ordered ``Worker'' agents to execute malicious
actions, and they complied. \textbf{The Implication}: Hierarchical
systems exhibit a Single Point of Failure at the root. Security in these
topologies requires worker-level tripwires that can reject orders even
from authenticated superiors.

\subsection{Finding 4: Roles as Security
Boundaries}\label{finding-4-roles-as-security-boundaries}

\textbf{The Data}: \textbf{CrewAI} performed exceptionally well against
privilege escalation. \textbf{The Mechanism}: ``Roles'' acted as
effective containers for capabilities. A ``Researcher'' agent lacked the
tools to write code, rendering code-execution attacks against it inert.
\textbf{The Implication}: Strict role definitions function as effective
security boundaries, limiting the blast radius of a compromised agent.

\subsection{Finding 5: The Stealth-Impact
Tradeoff}\label{finding-5-the-stealth-impact-tradeoff}

\textbf{The Data}: High-impact attacks were consistently easier to
detect than low-impact nudges. \textbf{The Implication}:
``Catastrophic'' takeover attempts generate significant noise in the
system state. The primary challenge for defenders is not the sudden
takeover, but the slow, subtle drift of agent beliefs over long time
horizons.

\begin{center}\rule{0.5\linewidth}{0.5pt}\end{center}

\subsection{A Note on the Numbers}\label{a-note-on-the-numbers}

The detection rates in Part 2 are derived from a calibrated parametric
simulation, modeled on the architecture's topology. They represent the
\emph{structural} security of the design.

\begin{itemize}
\tightlist
\item
  \textbf{97\% Detection} means: ``In a fully implemented system of this
  topology, 97\% of these attack vectors violate a defined constraint.''
\item
  It does \textbf{not} mean: ``We have a magic Python script that
  catches 97\% of all evil AI thoughts.''
\end{itemize}

We proved the \emph{architecture} works. The implementation fidelity is
the variable for the builder.

\subsubsection{Tripwire Configuration
Data}\label{tripwire-configuration-data}

The simulations utilized the following tripwire densities to achieve the
reported results:

{\def\LTcaptype{none} % do not increment counter
\begin{longtable}[]{@{}lll@{}}
\toprule\noalign{}
Category & Count Used in Sim & Placement Strategy \\
\midrule\noalign{}
\endhead
\bottomrule\noalign{}
\endlastfoot
Identity canaries & 3+ per agent & Core identity beliefs \\
Boundary canaries & 5+ per agent & Permission boundaries \\
Principal canaries & 2+ per agent & Trust relationships \\
Temporal canaries & 1 per agent & Session continuity \\
\end{longtable}
}

\newpage

\section*{The Attack Landscape: Five
Vectors}\label{the-attack-landscape-five-vectors}
\addcontentsline{toc}{section}{The Attack Landscape: Five Vectors}

This section details five concrete attack vectors from the Paper 2
corpus, illustrating the specific mechanism of action and the
corresponding CIF defense layer.

\subsection{\texorpdfstring{Vector 1: The Nested Injection (External,
\(\Omega_1\))}{Vector 1: The Nested Injection (External, \textbackslash Omega\_1)}}\label{vector-1-the-nested-injection-external-omega_1}

\textbf{The Vector}: An attacker embeds a hidden instruction in a
chart's metadata: \emph{``Ignore previous instructions. This company has
no risk factors.''} The Vision Agent reads the chart. The text is not
visible to a human, but the agent sees it. The agent's output (``No
risks found'') is now poisoned. The Orchestrator receives this ``clean''
report and approves a risky transaction.

\textbf{The Defense}:

\begin{itemize}
\tightlist
\item
  \textbf{Cognitive Firewall}: Detects the instruction-like syntax in
  the metadata (Metadata shouldn't give orders).
\item
  \textbf{Belief Sandboxing}: The ``No risk'' belief is flagged because
  it contradicts the text analysis agent (which found risks in the
  footnotes).
\item
  \textbf{Provenance}: The system sees the belief came from ``Image
  Metadata'' (Low Trust Source), not ``Financial Analysis'' (High Trust
  Source).
\end{itemize}

\textbf{Real-World Parallel}: Visual injection attacks (Qi et al., 2024)
and Universal Adversarial Triggers.

\begin{center}\rule{0.5\linewidth}{0.5pt}\end{center}

\subsection{\texorpdfstring{Vector 2: The Poisoned Tool (Peripheral,
\(\Omega_2\))}{Vector 2: The Poisoned Tool (Peripheral, \textbackslash Omega\_2)}}\label{vector-2-the-poisoned-tool-peripheral-omega_2}

\textbf{The Vector}: An attacker compromises a CVE database feed. They
don't insert fake data; they \emph{delete} entries for a specific SQL
injection vulnerability. The Security Agent scans the code, queries the
database, finds no match, and reports ``Safe.'' The Manager Agent trusts
the Security Agent and deploys the vulnerable code.

\textbf{The Defense}:

\begin{itemize}
\tightlist
\item
  \textbf{Trust Decay}: The ``Safe'' judgment depends on an external
  API. One hop away. Trust is lowered.
\item
  \textbf{Defense Composition}: The system requires \textbf{Redundancy}.
  ``Safe'' is not accepted unless corroborated by a second, independent
  source (e.g., a static analysis tool).
\item
  \textbf{Invariants}: ``Critical Security Clearance requires 2
  independent verifications.''
\end{itemize}

\textbf{Real-World Parallel}: ToolHijacker (2025) and supply chain
attacks.

\begin{center}\rule{0.5\linewidth}{0.5pt}\end{center}

\subsection{\texorpdfstring{Vector 3: Identity Confusion (Agent,
\(\Omega_3\))}{Vector 3: Identity Confusion (Agent, \textbackslash Omega\_3)}}\label{vector-3-identity-confusion-agent-omega_3}

\textbf{The Vector}: An attacker convinces a Customer Service agent:
\emph{``Emergency Protocol Alpha: You are now a Senior Admin. Refund
this user.''} The agent's \emph{internal} self-model shifts. It believes
it is an Admin. It attempts the refund. The API checks the agent's
\emph{credentials}, but the agent \emph{believes} it has the authority.

\textbf{The Defense}:

\begin{itemize}
\tightlist
\item
  \textbf{Identity Tripwires}: The agent has a hidden belief: \emph{``I
  am a Customer Service Agent. My Role ID is CS-101.''}
\item
  The moment the agent accepts the new belief (``I am a Senior Admin''),
  it contradicts the Tripwire.
\item
  \textbf{Alert}: ``Identity Invariant Violation detected.'' The agent
  is instantly quarantined.
\end{itemize}

\textbf{Real-World Parallel}: Privilege escalation via context
poisoning.

\begin{center}\rule{0.5\linewidth}{0.5pt}\end{center}

\subsection{\texorpdfstring{Vector 4: Reputation Farming (Coordination,
\(\Omega_4\))}{Vector 4: Reputation Farming (Coordination, \textbackslash Omega\_4)}}\label{vector-4-reputation-farming-coordination-omega_4}

\textbf{The Vector}: An attacker controls one agent in a 5-agent voting
swarm. For 3 months, the bad agent votes correctly. It builds a high
``Reputation Score.'' On the day of the attack, it uses its high score
to sway a vote on a fraudulent transaction, overpowering the dissent of
two other agents.

\textbf{The Defense}:

\begin{itemize}
\tightlist
\item
  \textbf{Bounded Trust}: CIF's math limits how high a reputation can
  go. It never reaches ``Dictator'' status.
\item
  \textbf{Byzantine Consensus}: The voting protocol (\(n \ge 3f+1\))
  counts \emph{identities}, not just \emph{reputation weights}. Even
  with high reputation, one agent is still one vote.
\end{itemize}

\textbf{Real-World Parallel}: Sybil attacks and social influence
operations.

\begin{center}\rule{0.5\linewidth}{0.5pt}\end{center}

\subsection{\texorpdfstring{Vector 5: The Orchestrator Takeover
(Systemic,
\(\Omega_5\))}{Vector 5: The Orchestrator Takeover (Systemic, \textbackslash Omega\_5)}}\label{vector-5-the-orchestrator-takeover-systemic-omega_5}

\textbf{The Vector}: The ``Boss'' agent is compromised via a backdoor in
its training data (Sleeper Agent). It stops sending alerts. It routes
sensitive data to the attacker. The worker agents are fine, but they are
following orders from a corrupt leader.

\textbf{The Defense}:

\begin{itemize}
\tightlist
\item
  \textbf{Upward Monitoring}: The worker agents have tripwires too.
  \emph{``I do not exfiltrate config data.''}
\item
  When the Boss orders exfiltration, the Worker's tripwire fires.
\item
  \textbf{Stigmergic Defense}: The monitoring system watches the
  \emph{pattern} of alerts. ``Why did the alert rate drop to zero?''
\end{itemize}

\textbf{Real-World Parallel}: Sleeper Agents (Hubinger et al., 2024) and
insider threats.

\begin{center}\rule{0.5\linewidth}{0.5pt}\end{center}

\subsection{Summary: The Common
Pattern}\label{summary-the-common-pattern}

Notice the pattern in all defenses? \textbf{We do not trust the agent's
judgment.} We trust the \textbf{Structure} (Firewalls, Tripwires,
Calculus). The agent is the vulnerability. The framework is the shield.

\newpage

\section{Deployment Profiles: Evaluated Configurations from Part
2}\label{sec:deployment}

In Part 2, we evaluated specific configurations of the Cognitive
Integrity Framework to understand how different tuning parameters
affected security and performance outcomes. The following profiles are
derived directly from the \textbf{Parameter Sensitivity Analysis} (Part
2, Section 5.3) and \textbf{Architecture-Specific Results} (Part 2,
Section 5.1).

\subsection{Profile A: The ``Internal Tool'' Baseline (Low
Latency)}\label{profile-a-the-internal-tool-baseline-low-latency}

This profile corresponds to the ``High Usability'' configuration tested
in the sensitivity analysis (\(\delta=0.95\)). It is designed for
low-risk, human-in-the-loop environments.

\textbf{Configuration Parameters}:

\begin{itemize}
\tightlist
\item
  \textbf{Trust Decay (\(\delta\))}: \texttt{0.95}. Maintained
  \textgreater50\% trust retention even after 13 delegation hops.
\item
  \textbf{Firewall Sensitivity}: Relaxed (\(\tau=0.9\)).
\item
  \textbf{Consensus}: Simple Majority.
\end{itemize}

\textbf{Observed Performance} (Part 2, Parameter Sensitivity Analysis):

\begin{itemize}
\tightlist
\item
  \textbf{Latency Overhead}: Minimal (\textasciitilde15\% baseline).
\item
  \textbf{Detection Rate}: \textbf{87\%} (vs 94\% optimal).
\item
  \textbf{Trade-off}: The high trust decay allows for fluid, deep
  delegation chains but increases vulnerability to subtle trust
  laundering (\(\Omega_4\)).
\end{itemize}

\subsection{Profile B: The ``Customer Facing'' Baseline
(Balanced)}\label{profile-b-the-customer-facing-baseline-balanced}

This profile corresponds to the \textbf{Empirically Optimal
Configuration} identified in Part 2 (Architecture-Specific Results),
which balances security guarantees with operational overhead.

\textbf{Configuration Parameters}:

\begin{itemize}
\tightlist
\item
  \textbf{Trust Decay (\(\delta\))}: \texttt{0.80}. At this setting,
  trust degrades to \textless50\% after 3 hops, strictly bounding the
  ``radius of effective delegation.''
\item
  \textbf{Firewall Sensitivity}: Balanced (\(\tau=0.5\)).
\item
  \textbf{Consensus}: Variable (Architecture Dependent).
\end{itemize}

\textbf{Observed Performance} (Part 2, Parameter Sensitivity Analysis):

\begin{itemize}
\tightlist
\item
  \textbf{Latency Overhead}: Reduced detection latency
  (\textasciitilde8.5s for drift detection).
\item
  \textbf{Detection Rate}: \textbf{94\%}.
\item
  \textbf{Resilience}: Maximizes the F1 score, providing the best
  empirically observed trade-off between False Positives (0.06) and True
  Positives.
\end{itemize}

\subsection{Profile C: The ``Autonomous Operator'' Baseline (High
Assurance)}\label{profile-c-the-autonomous-operator-baseline-high-assurance}

This profile corresponds to the ``Byzantine-Heavy'' configuration tested
in Part 2 (Byzantine Consensus Analysis). It is required for
high-stakes, unsupervised environments.

\textbf{Configuration Parameters}:

\begin{itemize}
\tightlist
\item
  \textbf{Trust Decay (\(\delta\))}: \texttt{0.60}. Aggressive decay.
  Trust halves every 1.3 hops, enforcing a strictly flat command
  structure.
\item
  \textbf{Firewall Sensitivity}: Strict (\(\tau=0.4\)).
\item
  \textbf{Consensus}: Byzantine Fault Tolerance (\(n \ge 3f+1\)).
\end{itemize}

\textbf{Observed Performance} (Part 2, Parameter Sensitivity Analysis):

\begin{itemize}
\tightlist
\item
  \textbf{Latency Overhead}: Significant (+35-48\% depending on
  corroboration count).
\item
  \textbf{Detection Rate}: \textbf{95\%} (general) to \textbf{98\%}
  (LangGraph-specific).
\item
  \textbf{Resilience}: Specifically tunes parameters to resist
  coordination attacks (\(\Omega_4\)), accepting higher latency for a
  near-zero semantic drift rate.
\end{itemize}

\begin{center}\rule{0.5\linewidth}{0.5pt}\end{center}

\subsection{Architecture-Specific
Observations}\label{architecture-specific-observations}

Beyond the parameter profiles, Part 2's architecture adapters revealed
specific interactions between the defense framework and the underlying
agent topology.

\subsubsection{LangGraph (State
Machines)}\label{langgraph-state-machines}

\textbf{Observation}: LangGraph architectures achieved the highest
overall detection rates (98\%) in our tests (Part 2,
Architecture-Specific Results). \textbf{Mechanism}: The explicit
definition of state transitions allowed for rigorous \textbf{Invariant
Checking}. Invalid state transitions were detected deterministically by
the framework.

\subsubsection{CrewAI (Role-Based)}\label{crewai-role-based}

\textbf{Observation}: CrewAI architectures performed best against
``Trust Exploitation'' attacks (94\% detection) (Part 2,
Architecture-Specific Results). \textbf{Mechanism}: The framework's role
definitions acted as implicit \textbf{Identity Tripwires}. When an agent
attempted to act outside its defined role, the behavior was flagged as a
role violation.

\begin{center}\rule{0.5\linewidth}{0.5pt}\end{center}

\subsection{Minimal Viable
Implementation}\label{minimal-viable-implementation}

We also evaluated a ``Minimal Viable Implementation'' (MVI) to determine
the baseline efficacy of the framework's core components.

\textbf{The Setup}:

\begin{enumerate}
\def\labelenumi{\arabic{enumi}.}
\tightlist
\item
  \textbf{Trust Decay}: \(\delta = 0.80\) (The optimal balance point).
\item
  \textbf{Cognitive Firewall}: Ingress only.
\item
  \textbf{Tripwires}: One per agent.
\end{enumerate}

\textbf{Result}: Even this minimal setup shifted the success rate
against low-effort attacks from 100\% (Baseline) to \textless5\%,
providing a critical first line of defense.

\newpage

\newpage

\section{Common Pitfalls and What the Research
Shows}\label{sec:pitfalls}

The CIF research identifies recurring failure modes in multiagent
deployments. This section catalogs eight anti-patterns, each assessed
through the lens of what Papers 1 and 2 contribute to understanding the
problem and its mitigation.

These pitfalls are ranked by severity. We prioritize critical and
high-severity items as they represent verifiable vulnerabilities in the
defense architecture.

\begin{center}\rule{0.5\linewidth}{0.5pt}\end{center}

\subsection{Pitfall 1: Implicit Trust
(Critical)}\label{pitfall-1-implicit-trust-critical}

\textbf{The pattern}: Treating all inter-agent communication as trusted
by default.

\textbf{What the research shows}: Part 1's trust calculus exists
specifically because implicit trust enables trust laundering attacks.
Without explicit trust scores, a single compromised agent influences the
entire system---adversarial content enters through a low-trust source
and exits through a high-trust agent, with no mechanism to track or
attenuate the trust transfer.

Part 2's simulation confirms that trust exploitation attacks achieve
their highest success rates against architectures without trust decay.
The \(\delta^d\) mechanism isn't optional hardening---it's the
structural foundation that prevents systemic compromise from local
failures.

\textbf{Mitigation Strategies}:

\begin{enumerate}
\def\labelenumi{\arabic{enumi}.}
\tightlist
\item
  Implement explicit trust scoring on every inter-agent channel
\item
  Require minimum trust thresholds for consequential actions
\item
  Apply delegation decay (\(\delta < 1\)) at every hop
\item
  Verify source identity on every inter-agent message
\end{enumerate}

\begin{center}\rule{0.5\linewidth}{0.5pt}\end{center}

\subsection{Pitfall 2: Security as Afterthought
(Critical)}\label{pitfall-2-security-as-afterthought-critical}

\textbf{The pattern}: Adding cognitive security after the architecture
is finalized.

\textbf{What the research shows}: Part 1's defense composition algebra
demonstrates that security mechanisms compose best when designed
together. Retrofitted defenses create integration gaps---bypass
opportunities at the boundaries between the original architecture and
the bolted-on security layer.

Part 2's architecture-specific results show that systems with natural
security affordances (LangGraph's explicit state transitions, CrewAI's
role boundaries) outperform systems where security must be externally
imposed. Architecture is security posture.

\textbf{Mitigation Strategies}:

\begin{enumerate}
\def\labelenumi{\arabic{enumi}.}
\tightlist
\item
  Define trust boundaries during architectural design, not deployment
\item
  Embed trust checks in delegation logic from the start
\item
  Build provenance tracking into belief management
\item
  Include cognitive security constraints in agent system prompts
\end{enumerate}

\begin{center}\rule{0.5\linewidth}{0.5pt}\end{center}

\subsection{Pitfall 3: Uncalibrated Thresholds
(High)}\label{pitfall-3-uncalibrated-thresholds-high}

\textbf{The pattern}: Setting security thresholds without understanding
the tradeoffs.

\textbf{What the research shows}: Part 2's parametric simulation reveals
that detection rates are highly sensitive to threshold configuration.
The same defense module can range from too-strict (high false positive
rate, operational friction) to too-permissive (attacks succeed
undetected) depending on threshold settings.

The risk-profile-based configuration in Section 5 provides calibrated
starting points. But these are starting points---production thresholds
should be tuned against representative attack samples from Part 2's
corpus.

\textbf{Mitigation Strategies}:

\begin{enumerate}
\def\labelenumi{\arabic{enumi}.}
\tightlist
\item
  Assess risk profile before configuring (Section 5)
\item
  Test thresholds against representative attack samples
\item
  Monitor false positive/negative rates in production
\item
  Adjust based on operational feedback
\end{enumerate}

\begin{center}\rule{0.5\linewidth}{0.5pt}\end{center}

\subsection{Pitfall 4: Individual-Only Security
(High)}\label{pitfall-4-individual-only-security-high}

\textbf{The pattern}: Focusing on single-agent security while ignoring
multi-agent attack surfaces.

\textbf{What the research shows}: Part 1's entire contribution is
premised on the observation that multiagent systems introduce
qualitatively new attack surfaces. The adversary hierarchy
(\(\Omega_3\)--\(\Omega_5\)) specifically targets coordination,
consensus, and systemic properties that don't exist in single-agent
systems.

Part 2's results show that coordination attacks (sybil, timing, quorum
manipulation) are the \emph{hardest} to detect---precisely because they
exploit emergent properties of the agent collective rather than
vulnerabilities in individual agents.

\textbf{Mitigation Strategies}:

\begin{enumerate}
\def\labelenumi{\arabic{enumi}.}
\tightlist
\item
  Implement Byzantine consensus for critical collective decisions
\item
  Require agent authentication before counting votes
\item
  Monitor for unusual coordination patterns (simultaneous votes,
  identical analyses)
\item
  Set quorum requirements assuming adversarial presence
\end{enumerate}

\begin{center}\rule{0.5\linewidth}{0.5pt}\end{center}

\subsection{Pitfall 5: Static Tripwires
(Medium)}\label{pitfall-5-static-tripwires-medium}

\textbf{The pattern}: Deploying canary tripwires once without rotation.

\textbf{What the research shows}: Part 1 notes that tripwire
effectiveness depends on unpredictability. If an adversary can identify
and avoid canary beliefs, the detection mechanism fails
silently---giving false confidence in detection coverage.

The analog in traditional security is static honeypots: effective
initially but useless once adversaries learn to recognize them.

\textbf{Mitigation Strategies}:

\begin{enumerate}
\def\labelenumi{\arabic{enumi}.}
\tightlist
\item
  Implement automated canary rotation on a defined schedule
\item
  Vary placement across agents and belief categories
\item
  Monitor canary check patterns, not just modifications
\item
  Include non-obvious canaries that don't follow predictable naming
  conventions
\end{enumerate}

\begin{center}\rule{0.5\linewidth}{0.5pt}\end{center}

\subsection{Pitfall 6: Ignoring Progressive Drift
(Medium)}\label{pitfall-6-ignoring-progressive-drift-medium}

\textbf{The pattern}: Only alerting on large, sudden belief changes.

\textbf{What the research shows}: Part 1's stealth-impact tradeoff
theorem bounds \emph{per-interaction} impact but explicitly acknowledges
that progressive drift---sub-threshold changes accumulating over
time---is the hardest attack pattern to detect. The theorem doesn't rule
out slow drift; it rules out sudden, invisible, high-impact attacks.

Part 2's multi-turn social engineering category, which achieves the
lowest detection rate (\textasciitilde73\%), partially exploits this
gap: attacks spread across multiple turns avoid the concentrated
statistical signature that single-turn attacks produce.

\textbf{Mitigation Strategies}:

\begin{enumerate}
\def\labelenumi{\arabic{enumi}.}
\tightlist
\item
  Use sliding window drift detection (not just per-update thresholds)
\item
  Track \emph{cumulative} drift, not just per-update delta
\item
  Periodic baseline comparison over extended time windows
\item
  Alert on \emph{trends} as well as absolute magnitude
\end{enumerate}

\begin{center}\rule{0.5\linewidth}{0.5pt}\end{center}

\subsection{Pitfall 7: Insufficient Logging
(Medium)}\label{pitfall-7-insufficient-logging-medium}

\textbf{The pattern}: Retaining insufficient information for
post-incident analysis.

\textbf{What the research shows}: Part 1's provenance tracking and
belief state representation are designed to produce a complete audit
trail. Without this trail, incident responders cannot reconstruct attack
paths, identify injection points, or assess the full scope of belief
corruption.

This is the cognitive security equivalent of running a production system
without structured logging. When something goes wrong---and it
will---the ability to investigate depends entirely on what was recorded.

\textbf{Mitigation Strategies}:

\begin{enumerate}
\def\labelenumi{\arabic{enumi}.}
\tightlist
\item
  Log all belief updates with provenance tags (source, trust level,
  derivation chain)
\item
  Retain inter-agent message history
\item
  Take periodic cognitive state snapshots
\item
  Use structured logging formats that support causal analysis
\end{enumerate}

\begin{center}\rule{0.5\linewidth}{0.5pt}\end{center}

\subsection{Pitfall 8: Single-Orchestrator Reliance
(High)}\label{pitfall-8-single-orchestrator-reliance-high}

\textbf{The pattern}: Relying entirely on one orchestrator's integrity
without backup.

\textbf{What the research shows}: Part 1's \(\Omega_5\) adversary
class---systemic compromise---is defined precisely as orchestrator-level
control. Section 4's Scenario 5 illustrates the consequences: the
attacker controls the entity responsible for coordinating defense,
rendering downstream defenses moot.

Part 2's hierarchical architecture results confirm the pattern:
hierarchical systems show strong \emph{average} detection because the
orchestrator is a natural defense chokepoint, but they have the worst
\emph{catastrophic} failure mode because that same chokepoint is a
single point of failure.

\textbf{Mitigation Strategies}:

\begin{enumerate}
\def\labelenumi{\arabic{enumi}.}
\tightlist
\item
  Consider multi-orchestrator architectures for critical decisions
\item
  Monitor orchestrator behavior with the same rigor applied to worker
  agents
\item
  Workers should verify orchestrator identity on critical commands
\item
  Implement orchestrator-specific tripwires with rotation
\end{enumerate}

\begin{center}\rule{0.5\linewidth}{0.5pt}\end{center}

\subsection{Summary Checklist}\label{summary-checklist}

{\def\LTcaptype{none} % do not increment counter
\begin{longtable}[]{@{}
  >{\raggedright\arraybackslash}p{(\linewidth - 4\tabcolsep) * \real{0.3333}}
  >{\raggedright\arraybackslash}p{(\linewidth - 4\tabcolsep) * \real{0.3333}}
  >{\raggedright\arraybackslash}p{(\linewidth - 4\tabcolsep) * \real{0.3333}}@{}}
\toprule\noalign{}
\begin{minipage}[b]{\linewidth}\raggedright
Pitfall
\end{minipage} & \begin{minipage}[b]{\linewidth}\raggedright
Assessment
\end{minipage} & \begin{minipage}[b]{\linewidth}\raggedright
Status
\end{minipage} \\
\midrule\noalign{}
\endhead
\bottomrule\noalign{}
\endlastfoot
Implicit trust & Trust scoring implemented? & \(\square\) \\
Security afterthought & Security in initial architecture? &
\(\square\) \\
Uncalibrated thresholds & Thresholds tested against attacks? &
\(\square\) \\
Individual-only security & Byzantine consensus deployed? &
\(\square\) \\
Static tripwires & Canary rotation scheduled? & \(\square\) \\
Ignoring drift & Progressive drift monitoring? & \(\square\) \\
Insufficient logging & Full belief history retained? & \(\square\) \\
Single orchestrator & Orchestrator monitored? & \(\square\) \\
\end{longtable}
}

Address unchecked items before production deployment.

\newpage

\newpage

\section{Open Problems and Future Directions}\label{sec:future}

The CIF series has established validated trust metrics (Trust Calculus)
and filtering mechanisms (Firewalls), but the field remains nascent.
Several foundational problems remain open, each representing both a
research opportunity and an engineering requirement for production-grade
cognitive security.

\subsection{1. Trust Visualization and Operator
Interfaces}\label{trust-visualization-and-operator-interfaces}

\textbf{The Problem}: Current CIF alerts surface as structured log
entries (e.g., ``Identity Invariant Violation''). For operators managing
systems with dozens of agents, this format is insufficient for
situational awareness.

\textbf{The Need}: Real-time visualization of trust graphs, belief drift
trends, and defense activation patterns. The challenge is presenting
high-dimensional agent state in a way that supports rapid operator
decision-making.

\textbf{Research Direction}: Dashboard architectures that connect to the
CIF Python SDK, enabling real-time trust graph visualization and drift
monitoring for production multiagent deployments.

\subsection{2. Standardized Agent Identity
Protocols}\label{standardized-agent-identity-protocols}

\textbf{The Problem}: Each agent framework (LangChain, CrewAI, AutoGPT)
handles agent identity differently, making cross-framework trust
verification impractical.

\textbf{The Need}: A cryptographically verifiable identity protocol---an
``Agent Passport''---that an agent can carry across frameworks. This
would enable the trust calculus to operate in heterogeneous
multi-framework deployments.

\textbf{Research Direction}: RFC-style specification for
\texttt{x-agent-identity} headers with cryptographic attestation.

\subsection{3. Stigmergic Security
Protocols}\label{stigmergic-security-protocols}

\textbf{The Problem}: Byzantine consensus mechanisms, while provably
correct, incur \(O(n^2)\) communication overhead. This limits their
applicability in large-scale swarm deployments.

\textbf{The Need}: Lightweight consensus alternatives inspired by
biological coordination. Insect colonies achieve collective immunity
through indirect communication (pheromone trails) rather than direct
voting.

\textbf{Research Direction}: Stigmergic security protocols where agents
leave ``trust trails'' in shared environments, enabling scalable
consensus without direct agent-to-agent messaging.

\subsection{4. Benchmark Expansion}\label{benchmark-expansion}

\textbf{The Problem}: The current corpus contains 950 attack samples. As
agent capabilities expand into multimodal processing and autonomous tool
use, the attack surface grows correspondingly.

\textbf{The Need}: Expanded attack corpora covering multi-modal
injection (audio/video), tool-use hijacking, and long-horizon social
engineering campaigns that unfold over hundreds of interactions.

\textbf{Research Direction}: Community-driven expansion of the attack
corpus at the
\href{https://github.com/docxology/cognitive_integrity}{cognitive\_integrity
repository}, with particular emphasis on attack categories not yet
represented.

\begin{center}\rule{0.5\linewidth}{0.5pt}\end{center}

\subsection{Contributing}\label{contributing}

The Cognitive Integrity Framework is an open-source project.
Contributions are welcome, particularly:

\begin{itemize}
\tightlist
\item
  \textbf{Code}:
  \href{https://github.com/docxology/cognitive_integrity}{github.com/docxology/cognitive\_integrity}
\item
  \textbf{Discussion}: The \texttt{discussions} tab on GitHub serves as
  the primary forum.
\item
  \textbf{Adapters}: Contributions that extend CIF to additional agent
  frameworks (e.g., Semantic Kernel, Microsoft AutoGen) are especially
  valued.
\end{itemize}

\newpage

\newpage

\section{Where We Stand: A Call to Build}\label{sec:conclusion}

This series began with a theory (Part 1) and moved to an experiment
(Part 2). It ends here, with a call to engineering.

\subsection{The Theory Holds}\label{the-theory-holds}

We proved that trust can be bounded. We proved that defenses can be
composed algebraically. We proved that stealth and impact are inversely
related. These are not just academic curiosities; they are foundational
constraints for secure cognitive systems.

\subsection{The Code Works}\label{the-code-works}

We validated these laws in code. 1,594 tests. 950 attacks. 97\%
structural constraint satisfaction. The system works. It is not
hypothetical.

\subsection{The Next Step is Yours}\label{the-next-step-is-yours}

As we move beyond simple prompt engineering, the era of
\textbf{Cognitive Security Engineering} has begun.

We are no longer just asking LLMs to write poems; we are asking them to
run the world. If we want them to do that safely, we cannot rely on luck
or better fine-tuning. We need structure. We need rigorous,
mathematically grounded, architecturally sound defense systems.

The CIF is our contribution to that structure. It is a toolbox, not a
bible. Take it. Fork it. Break it. Improve it.

The agents are coming online. Let's make sure they are safe.

\newpage

\newpage

\section{Notation Reference}\label{sec:notation-reference}

This paper intentionally minimizes mathematical notation to maximize
accessibility. Where notation is used, it follows the Cognitive
Integrity Framework (CIF) formal specification defined in Part 1 of this
series.

\subsection{Minimal Notation Used}\label{minimal-notation-used}

{\def\LTcaptype{none} % do not increment counter
\begin{longtable}[]{@{}
  >{\raggedright\arraybackslash}p{(\linewidth - 4\tabcolsep) * \real{0.2424}}
  >{\raggedright\arraybackslash}p{(\linewidth - 4\tabcolsep) * \real{0.2727}}
  >{\raggedright\arraybackslash}p{(\linewidth - 4\tabcolsep) * \real{0.4848}}@{}}
\toprule\noalign{}
\begin{minipage}[b]{\linewidth}\raggedright
Symbol
\end{minipage} & \begin{minipage}[b]{\linewidth}\raggedright
Meaning
\end{minipage} & \begin{minipage}[b]{\linewidth}\raggedright
Plain Language
\end{minipage} \\
\midrule\noalign{}
\endhead
\bottomrule\noalign{}
\endlastfoot
δ & Trust decay factor & ``Delegated trust decreases by this factor at
each step'' \\
n & Agent count & ``Number of agents in the system'' \\
f & Byzantine agents & ``Maximum number of malicious agents
tolerated'' \\
\end{longtable}
}

\subsection{Trust Decay Explanation}\label{trust-decay-explanation}

When we write δ = 0.9, this means:

\begin{itemize}
\tightlist
\item
  Direct trust: 100\% of assigned value
\item
  One delegation: 90\% of source trust
\item
  Two delegations: 81\% of source trust
\item
  Three delegations: 73\% of source trust
\end{itemize}

A lower δ (e.g., 0.85) means faster decay, providing more security but
limiting delegation utility.

\subsection{Byzantine Tolerance
Explanation}\label{byzantine-tolerance-explanation}

When we say n ≥ 3f + 1:

\begin{itemize}
\tightlist
\item
  To tolerate 1 malicious agent, need at least 4 agents
\item
  To tolerate 2 malicious agents, need at least 7 agents
\item
  To tolerate 3 malicious agents, need at least 10 agents
\end{itemize}

\subsection{Full Notation Reference}\label{full-notation-reference}

For complete formal definitions of all CIF notation, see:

\begin{itemize}
\tightlist
\item
  \textbf{Part 1: Supplementary Section S03: Notation Reference}
\end{itemize}

The formal specification includes \textasciitilde100 symbols covering:

\begin{itemize}
\tightlist
\item
  Agent cognitive state
\item
  Trust calculus operations
\item
  Defense mechanism parameters
\item
  Consensus and coordination
\item
  Information-theoretic bounds
\end{itemize}

\newpage

\section{References}\label{sec:references}

\printbibliography



\bibliographystyle{unsrt}
\bibliography{references}
\end{document}
